% JFM Rapids draft: A spectral theory for turbulent burning velocity across scales
% Built from the authors' JFM template (front and back matter as supplied).
% Figures: figures/ (all generated by make_figures.ipynb).
% Open items for the authors are marked "% TODO".
% Class: jfm.cls + jfm.bst from the authors' (Overleaf) template, as in JFM-template.tex. The newer class in
% ../JFM-FLM-AUTemplate.zip (JFM-FLM_Au.cls) has no \backsection; switching needs \corresau and bmhead instead.

\documentclass[lineno]{jfm}

\usepackage{graphicx}
%\usepackage{epstopdf,epsfig}
\usepackage{newtxtext}
\usepackage{newtxmath}
\usepackage{natbib}
\defcitealias{Tamadonfar2014}{Tamadonfar \& G. 2014} % short form for table 1
\usepackage{hyperref}
\hypersetup{
    colorlinks = true,
    urlcolor   = blue,
    citecolor  = black,
}
\newtheorem{lemma}{Lemma}
\newtheorem{corollary}{Corollary}

% notation
\providecommand{\Rey}{\mbox{\textit{Re}}}
\providecommand{\Ka}{\mbox{\textit{Ka}}}
\newcommand{\Sct}{\mbox{\textit{Sc}}_{\, T}}
\newcommand{\sL}{s_L}
\newcommand{\dL}{\delta_L}
\newcommand{\Cs}{\mathcal{C}}
\newcommand{\kh}{\hat{\kappa}}
\newcommand{\De}{D_e^{+}}
\newcommand{\up}{u^{\prime}}
\newcommand{\sym}[1]{\raisebox{-0.55ex}{\includegraphics[height=2.4ex]{sym_#1}}} % plot symbol, figures/sym_*.pdf

\title{A spectral theory for turbulent burning velocity across scales}

\author{Michael D. Walker\aff{1,2}
  \corresp{\email{michael.walker@princeton.edu}},
  G. V. Nivarti\aff{2},
  M. E. Mueller\aff{1},
  R. S. Cant\aff{2},
 \and S. Hochgreb\aff{2}}

\affiliation{\aff{1}Department of Mechanical and Aerospace Engineering, Princeton University, Princeton, NJ 08544, USA
\aff{2}Department of Engineering, University of Cambridge, Trumpington Street, Cambridge CB2 1PZ, UK}

\begin{document}
\maketitle

\begin{abstract}
Damk\"ohler's two hypotheses attribute the enhancement of the turbulent burning velocity $s_T$ over the laminar value $\sL$ to an increase in flame-surface area by eddies larger than the flame thickness $\dL$, and to an increase in effective diffusivity by smaller eddies. In intense turbulence, the measured $s_T / \sL$ exceeds the measured area ratio several-fold, an excess that Nivarti \textit{et al.} (\textit{J.~Fluid Mech.}, vol.~858, 2019, R1) attributed to a sub-$\dL$ diffusivity integrated over a model energy spectrum. Here that integral is extended to the whole spectrum and made dimensionally consistent by taking the diffusivity at each wavenumber from the spectral eddy diffusivity of two-point closures. For a model spectrum, the total enhancement reduces to $\Cs \, \mathcal{F}(\Rey) \, \up / \sL$, where $\up$ is the turbulence intensity, $\Cs$ is a constant and $\mathcal{F}$ is a weak, saturating function of the turbulence Reynolds number $\Rey$. The resulting relation $s_T / \sL = [1 + \Cs \, \mathcal{F}(\Rey) \, \up / \sL]^{1/2}$ recovers Damk\"ohler's linear growth of $s_T$ with $\up$ in weak turbulence and the sublinear `bending' in strong turbulence, and its sub-$\dL$ part recovers the $\Ka^{1/3}$ scaling of Nivarti \textit{et al.}, where $\Ka$ is the Karlovitz number. The relation has the same form as the power-law flame-wrinkling factor of Charlette \textit{et al.} (\textit{Combust. Flame}, vol.~131, 2002, p.~159) with exponent $1/2$, which supports the square-root dependence at large scales. With spectral constants from non-reacting turbulence, the prediction follows 80 measurements from eight datasets spanning four decades of $\Ka$ and $\Rey$, 71\,\% of them within a factor of 1.5.
\end{abstract}

\begin{keywords}
combustion, turbulent reacting flows, turbulent mixing
\end{keywords}

% {\bf MSC Codes }  {\it(Optional)} Please enter your MSC Codes here

%=====================================================================
\section{Introduction}\label{sec:intro}

\citet{Damkohler1940} proposed two mechanisms by which turbulence raises the burning velocity of a premixed flame above its laminar value $\sL$. Eddies larger than the laminar flame thickness $\dL$ wrinkle the flame without altering its inner structure, so that the turbulent burning velocity $s_T$ grows with the flame-surface area,
\begin{equation}
  \frac{s_T}{\sL} = \frac{A_T}{A} = 1 + \frac{\Delta A}{A},
  \label{eq:DH1}
\end{equation}
where $A_T$ is the wrinkled area and $A$ the mean (unwrinkled) area. Eddies smaller than $\dL$ instead enhance transport within the flame as a diffusive process. Since $\sL \propto D^{1/2}$ for a fixed reaction time \citep{ZeldovichFK1938}, an effective turbulent diffusivity $D_T = D + \Delta D$ yields
\begin{equation}
  \frac{s_T}{\sL} = \left(\frac{D_T}{D}\right)^{1/2} = \left(1 + \frac{\Delta D}{D}\right)^{1/2}.
  \label{eq:DH2}
\end{equation}
The two hypotheses, (\ref{eq:DH1}) and (\ref{eq:DH2}), are conventionally assigned to separate regimes, delineated by the Karlovitz number $\Ka \equiv \dL^2 / \eta^2$, where $\eta$ is the Kolmogorov scale \citep{Peters1999,Peters2000}. Throughout, unity Schmidt and Lewis numbers are assumed, so that $\nu = D = \sL \dL$, and the Reynolds and Karlovitz numbers are
\begin{equation}
  \Rey \equiv \frac{\up \ell_0}{\nu} = \frac{\up}{\sL} \frac{\ell_0}{\dL}, \quad
  \Ka = \left(\frac{\up}{\sL}\right)^{3/2} \left(\frac{\ell_0}{\dL}\right)^{-1/2},
  \quad \mbox{such that} \quad \frac{\up}{\sL} = \Ka^{1/2} \Rey^{1/4},
  \label{eq:groups}
\end{equation}
with $\up$ the r.m.s.\ velocity fluctuation and $\ell_0$ the integral length scale.

Planar measurements in intense turbulence \citep{YuenGulder2013,Wabel2017} show that (\ref{eq:DH1}) fails for $\Ka > 1$: the area ratio $A_T / A$ saturates near 4--6 while $s_T / \sL$ continues to rise, reaching nearly 30 (figure~\ref{fig:data}\textit{b}; see also the reviews of \citealt{Driscoll2020} and \citealt{Hochgreb2023}). Direct numerical simulations at $\up / \sL = 35$ show a smaller excess, about 14\,\% on average, attributed to flame collisions \citep{PoludnenkoOran2011,Hicks2015}. \citet{Nivarti2019} reconciled  $s_T$ and $A_T/A$ by writing $s_T / \sL = (A_T / A)(1 + \Delta D^{>} / D)^{1/2}$ and estimating the sub-$\dL$ enhancement $\Delta D^{>}$ by summing Prandtl mixing-length diffusivities over the portion of a model energy spectrum with wavenumbers $\kappa > \dL^{-1}$. The result, $\Delta D^{>} / D \sim \Ka^{2/3}$, accounted for the measured ratio $(s_T / \sL) / (A_T / A) \sim \Ka^{1/3}$. Related spectral estimates of small-scale enhancement have been made by \citet{Zimont1979}, \citet{RonneyYakhot1992} and \citet{Gulder2007}, and spectral closures have previously been used for the large-scale regime \citep{Peters1992,Chaudhuri2011,Chaudhuri2012}.

Because area enhancement by large eddies is itself a function of the spectrum, it is natural to ask whether a single spectral sum over all scales can account for $s_T$ in every regime \citep{Walker2019}. Two difficulties must be overcome. First, the scale velocity $u(\kappa) = [2 E(\kappa)]^{1/2}$ used by \citet{Nivarti2019}, where $E(\kappa)$ is the energy spectrum function, has dimensions of $\mbox{L}^{3/2} \, \mbox{T}^{-1}$, so the integrand $u(\kappa) / \kappa$ is not a diffusivity and its prefactor is arbitrary. Second, the square root in (\ref{eq:DH2}) derives from laminar-flame theory, and there is no obvious reason for it to hold at scales where area enhancement contributes significantly to the burning velocity.

This paper resolves both. In \S\,\ref{sec:formulation} the scale-dependent diffusivity is taken from the spectral eddy diffusivity of two-point closures \citep{Kraichnan1976,CholletLesieur1981,Chollet1984}, which fixes its dimensions and prefactor. In \S\,\ref{sec:analysis} the resulting integral is shown to reduce to a function of $\up / \sL$ and $\Rey$ alone, and the reduced relation is shown to coincide with the power-law flame-wrinkling factor of \citet{Charlette2002a} with exponent $1/2$. Section~\ref{sec:comparison} compares the prediction with eight experimental datasets.

\begin{figure}
  \centering
  \begin{minipage}[t]{0.50\textwidth}\raggedright (\textit{a})\\[-2pt]
    \includegraphics[width=\linewidth]{fig1a_regime}\end{minipage}\hfill
  \begin{minipage}[t]{0.49\textwidth}\raggedright (\textit{b})\\[-2pt]
    \includegraphics[width=\linewidth]{fig1b_area}\end{minipage}
  \caption{(\textit{a}) Experimental conditions on the turbulent combustion regime diagram. (\textit{b}) Burning-velocity ratio $s_T / \sL$ (black) and flame-surface-area ratio $A_T / A$ (blue, with $\pm 20 \, \%$ uncertainty from \citealt{Wabel2017thesis,Yuen2009,Filatyev2005}) against $\Ka$ for the datasets where both were measured; for \citet{Filatyev2005}, $A_T / A$ is the squared ratio of wrinkled to mean flame perimeter, multiplied by a $2 / \pi$ consistency factor for comparing slot and round burners. Symbols and data sources are listed in table~\ref{tab:data}.}
  \label{fig:data}
\end{figure}

%=====================================================================
\section{Formulation}\label{sec:formulation}

\subsection{Spectral eddy diffusivity}

In two-point closures of isotropic turbulence, the depletion of energy from a mode $\kappa$ to all modes beyond a cut-off $\kappa_c > \kappa$ is represented by a spectral eddy viscosity \citep{Kraichnan1976,CholletLesieur1981,Pope2000},
\begin{equation}
  \nu_e(\kappa | \kappa_c) = \nu_e^{+}(\kappa / \kappa_c) \left[\frac{E(\kappa_c)}{\kappa_c}\right]^{1/2},
  \label{eq:nue}
\end{equation}
where the non-dimensional function $\nu_e^{+}$, computed with the eddy-damped quasi-normal Markovian closure, is flat for $\kappa \ll \kappa_c$ and rises to a cusp as $\kappa \to \kappa_c$ \citep{Chollet1984}. \citet{Chollet1984} also computed the corresponding eddy diffusivity $\De$ of a passive scalar, whose ratio to $\nu_e^{+}$ defines a spectral turbulent Schmidt number $\Sct = \nu_e^{+} / \De$. Fits to these curves (figure~\ref{fig:spectrum}\textit{b}; coefficients in supplementary table~S1) give, at the cusp, $\nu_e^{+}(1) = 0.868$, $\Sct(1) = 0.462$ and hence $\De(1) = 1.88$.
% TODO: these cusp values come from the authors' fit to digitised Chollet (1984) curves; Pope's regression,
% nu_e^+ = C^{-3/2}[0.441 + 15.2 exp(-3.03 kappa_c/kappa)], gives nu_e^+(1) = 0.640. State which source is
% definitive and give the fit coefficients (Spectral_Viscosity.ipynb, cell 4) in the supplementary material.

The factor $[E(\kappa) / \kappa]^{1/2}$ in (\ref{eq:nue}) can be written $u_\kappa / \kappa$, where $u_\kappa \equiv [\kappa E(\kappa)]^{1/2}$ is the characteristic velocity of eddies of size $\kappa^{-1}$. It is therefore Prandtl's mixing-length product of a velocity and a length at that scale, with the dimensions of a diffusivity. Accordingly, we assign to each wavenumber the scalar eddy diffusivity at the cusp,
\begin{equation}
  D_\kappa(\kappa) \equiv \De(1) \left[\frac{E(\kappa)}{\kappa}\right]^{1/2} = \De(1) \, \frac{u_\kappa}{\kappa},
  \label{eq:Dk}
\end{equation}
that is, the transport exerted by eddies of wavenumber $\kappa$ on scales just larger than themselves. Equation (\ref{eq:Dk}) replaces the integrand $[2 E(\kappa)]^{1/2} / \kappa$ of \citet{Nivarti2019}, removing the dimensional inconsistency and taking the prefactor from the closure rather than setting it to unity.

\subsection{Total diffusivity enhancement}

Following \citet{Nivarti2019}, the contributions of all scales are summed with wavenumber measured in units of the flame wavenumber $\dL^{-1}$. Using $D = \sL \dL$,
\begin{equation}
  \frac{\Delta D}{D} = \frac{1}{D} \int_{\kappa_0}^{\kappa_\eta} D_\kappa \, \mathrm{d}(\kappa \dL)
  = \frac{\De(1)}{\sL} \int_{\kappa_0}^{\kappa_\eta} u_\kappa \, \mathrm{d} \ln \kappa,
  \label{eq:DD}
\end{equation}
with $\kappa_0 = \ell_0^{-1}$ and $\kappa_\eta = \eta^{-1}$. This weighting is the central modelling assumption of the theory: each logarithmic band of eddies adds a diffusivity $\De \, u_\kappa \dL$, with the flame thickness as its mixing length. It differs from the classical estimate $\Delta D \propto \up \ell_0$, in which every eddy mixes over its own size \citep{Damkohler1940}. The two can be compared through the flux each band carries across the flame. Let $\Delta \Lambda$ be the change in a progress variable $\Lambda$ across the flame, so that the molecular flux is $D \Delta \Lambda / \dL = \sL \Delta \Lambda$. A band of velocity $u_\kappa$ can transport at most $u_\kappa \Delta \Lambda$. With mixing length $\kappa^{-1}$, its gradient-diffusion flux would be $u_\kappa(\kappa \dL)^{-1} \Delta \Lambda$. For eddies smaller than the flame ($\kappa \dL > 1$) this is below that bound, and replacing $\kappa^{-1}$ by $\dL$ changes only an $\mathcal{O}(1)$ constant, because both sums are dominated by $\kappa \approx \dL^{-1}$. For eddies larger than the flame the gradient-diffusion flux exceeds the bound, so their mixing length is capped at $\dL$. Each band then carries a flux of order $u_\kappa \Delta \Lambda$, which relative to the molecular flux gives (\ref{eq:DD}). The estimate $\Delta D \propto \up \ell_0$ is therefore valid only in Damk\"ohler's limit $\ell_0 < \dL$. A further consequence of (\ref{eq:DD}) is that $\dL$ is absent from the enhancement factor, which becomes a ratio of velocities.

For the energy spectrum, we adopt the model of \citet{Pope2000} used by \citet{Nivarti2019},
\begin{equation}
  E(\kappa) = C \varepsilon^{2/3} \kappa^{-5/3} f_L(\kappa \ell_0) f_\eta(\kappa \eta), \quad
  f_L = \left[1 + \frac{c_L}{(\kappa \ell_0)^{2}}\right]^{-(5/3+p_0)/2}, \quad
  f_\eta = \mathrm{e}^{-\beta \kappa \eta},
  \label{eq:spectrum}
\end{equation}
where $\varepsilon$ is the kinetic energy dissipation rate. In units of the laminar flame, with $\kh \equiv \kappa \dL$, the dissipation rate $\varepsilon = 15 \nu(\up / \lambda)^2$ with Taylor microscale $\lambda = \ell_0(10 / \Rey)^{1/2}$ \citep{Taylor1935} becomes $\hat \varepsilon \equiv \varepsilon \dL / \sL^3 = (3/2) \Ka^2$ \citep{Nivarti2019}. The integration limits become $\kh_0 = \Ka^{1/2} \Rey^{-3/4}$ and $\kh_\eta = \Ka^{1/2}$ \citep{Walker2019}. Substituting (\ref{eq:spectrum}) into (\ref{eq:DD}) gives
\begin{equation}
  \mathcal{S} \equiv \frac{\Delta D}{D} = \Cs \, \Ka^{2/3} \int_{\Ka^{1/2} \Rey^{-3/4}}^{\Ka^{1/2}}
  \kh^{-4/3} \, \hat f_L^{1/2} \hat f_\eta^{1/2} \, \mathrm{d} \kh,
  \qquad \Cs \equiv \De(1) \, C^{1/2}(3/2)^{1/3},
  \label{eq:S}
\end{equation}
where $\hat f_L = [1 + c_L \Ka \Rey^{-3/2} \kh^{-2}]^{-(5/3+p_0)/2}$ and $\hat f_\eta = \exp(-\beta \kh \Ka^{-1/2})$. Compared with \citet[][their equation (2.7)]{Nivarti2019}, the exponent of the integrand changes from $-11/6$ to $-4/3$, and the prefactor $(2 C)^{1/2}(3/2)^{1/3}$ is replaced by $\Cs$, which equals 2.64 for $C = 1.5$. Finally, (\ref{eq:DH2}) is applied over the entire spectrum,
\begin{equation}
  \frac{s_T}{\sL} = \left(1 + \mathcal{S}\right)^{1/2},
  \label{eq:sT}
\end{equation}
a step justified in \S\,\ref{sec:wrinkling}.

\begin{figure}
  \centering
  \begin{minipage}[t]{0.49\textwidth}\raggedright (\textit{a})\\[-2pt]
    \includegraphics[width=\linewidth]{fig2a_spectrum}\end{minipage}\hfill
  \begin{minipage}[t]{0.49\textwidth}\raggedright (\textit{b})\\[-2pt]
    \includegraphics[width=\linewidth]{fig2b_chollet}\end{minipage}
  \caption{(\textit{a}) Model energy spectrum (\ref{eq:spectrum}) in flame units for increasing $p_0$. Wavenumbers below $\dL^{-1}$ (red) contribute a mixed $\Rey$--$\Ka$ dependence to (\ref{eq:S}), and those above $\dL^{-1}$ (blue) a dependence on $\Ka$ alone. (\textit{b}) Spectral eddy viscosity $\nu_e^{+}$, eddy diffusivity $\De$ and turbulent Schmidt number $\Sct$ computed by \citet{Chollet1984}, the regression for $\nu_e^{+}$ given by \citet{Pope2000}, and those used here.}
  \label{fig:spectrum}
\end{figure}

%=====================================================================
\section{Analysis}\label{sec:analysis}

\subsection{Reduction to two groups}\label{sec:reduction}

Equation (\ref{eq:S}) simplifies considerably. Substituting $\kh = \kh_0 \kappa^*$, where $\kappa^* \equiv \kappa \ell_0$, and using $\Ka^{1/2} \Rey^{1/4} = \up / \sL$ from (\ref{eq:groups}), the prefactor $\Ka^{2/3} \kh_0^{-1/3}$ reduces to $\up / \sL$ and the integration limits become $1$ and $\Rey^{3/4}$. The enhancement is therefore
\begin{equation}
  \mathcal{S} = \Cs \, \frac{\up}{\sL} \, \mathcal{F}(\Rey), \qquad
  \mathcal{F}(\Rey) = \int_1^{\Rey^{3/4}} \kappa^{*-4/3} \left(1 + \frac{c_L}{\kappa^{*2}}\right)^{-(5/3+p_0)/4}
  \exp \left(-\frac{\beta \kappa^*}{2 \Rey^{3/4}}\right) \mathrm{d} \kappa^*,
  \label{eq:SI}
\end{equation}
and (\ref{eq:sT}) becomes
\begin{equation}
  \frac{s_T}{\sL} = \left[1 + \Cs \, \mathcal{F}(\Rey) \, \frac{\up}{\sL}\right]^{1/2}.
  \label{eq:main}
\end{equation}
This is equivalent to $\Delta D = \Cs \, \mathcal{F}(\Rey) \, \up \dL$, an eddy diffusivity with the r.m.s.\ velocity as its velocity scale and the flame thickness as its mixing length. The burning velocity thus depends on two groups only, $\up / \sL$ and $\Rey$; the Karlovitz number is not an independent parameter.

The function $\mathcal{F}$ carries the spectral contribution. With $f_L = f_\eta = 1$ (inertial range only), $\mathcal{F} = 3(1 - \Rey^{-1/4})$. Retaining the dissipation range ($f_L = 1$), $\mathcal{F}$ can be written in terms of the upper incomplete gamma function $\Gamma(s, x)$,
\begin{align}
  \mathcal{F} &= (\beta / 2)^{1/3} \Rey^{-1/4} \left[\Gamma(-1/3, \, \beta \Rey^{-3/4} / 2) - \Gamma(-1/3, \, \beta / 2)\right] \nonumber \\
  &= 3 - c_\beta \Rey^{-1/4} + \mathcal{O}(\Rey^{-3/4}),
  \label{eq:Igamma}
\end{align}
with $c_\beta = (\beta / 2)^{1/3}[\Gamma(-1/3, \, \beta / 2) - \Gamma(-1/3)] = 5.61$ for $\beta = 5.2$. The energy-containing range reduces the asymptote to $\mathcal{F}_\infty = \int_1^\infty \kappa^{*-4/3}(1 + c_L \kappa^{*-2})^{-(5/3+p_0)/4} \, \mathrm{d} \kappa^*$, which equals 2.56, 2.24 and 2.05 for $p_0 = 0$, 2 and 4 ($c_L = 6.78$). For $p_0 = 0$, $\mathcal{F} = 0.98$, 1.60 and 2.00 at $\Rey = 10^2$, $10^3$ and $10^4$. Since $\mathcal{F}$ approaches its asymptote as $\Rey^{-1/4}$, the Reynolds-number dependence persists up to the highest laboratory Reynolds numbers, but it is weak (figure~\ref{fig:comparison}\textit{a}). The dissipation range lowers $\mathcal{F}$ more than the energy-containing range for $\Rey \lesssim 10^3$, and less for larger $\Rey$.

\begin{figure}
  \centering
  \begin{minipage}[t]{0.50\textwidth}\raggedright (\textit{a})\\[-2pt]
    \includegraphics[width=\linewidth]{fig3a_F}\end{minipage}\hfill
  \begin{minipage}[t]{0.49\textwidth}\raggedright (\textit{b})\\[-2pt]
    \includegraphics[width=\linewidth]{fig3b_comparison}\end{minipage}
  \caption{(\textit{a}) The function $\mathcal{F}(\Rey)$ of (\ref{eq:SI}) for $p_0 = 0$, 2 and 4, with the inertial-range result $3(1 - \Rey^{-1/4})$ and the form (\ref{eq:Igamma}) with $f_L = 1$; short lines at the right mark $\mathcal{F}_\infty$, and shading spans the $\Rey$ of the measurements. (\textit{b}) Measured $s_T / \sL$ (black) and $A_T / A$ (blue) against the prediction (\ref{eq:main}). Dashed line, equality; dotted lines, a factor of 1.5 either side. Symbols as in table~\ref{tab:data}.}
  \label{fig:comparison}
\end{figure}

Three limits of (\ref{eq:main}) connect it with earlier results.

(i) \textit{Weak turbulence}, $\Cs \, \mathcal{F} \, \up / \sL \ll 1$: $s_T / \sL \approx 1 + \thalf \Cs \, \mathcal{F} \, \up / \sL$. This is Damk\"ohler's linear scaling $s_T - \sL \propto \up$, which he derived from (\ref{eq:DH1}) by taking the area increase proportional to $\up / \sL$.

(ii) \textit{Strong turbulence}, $\Cs \, \mathcal{F} \, \up / \sL \gg 1$: $s_T / \sL \approx (\Cs \, \mathcal{F} \, \up / \sL)^{1/2}$, i.e.\ $s_T / \sL \propto \Ka^{1/4} \Rey^{1/8}$ at fixed $\mathcal{F}$. The transition between (i) and (ii) is the `bending' of $s_T$ with increasing $\up$ \citep{Bradley1992,LipatnikovChomiak2002}. It arises here without any quenching mechanism, from the square root in (\ref{eq:sT}) alone. For comparison, taking $\Delta D \propto \up \ell_0$ in (\ref{eq:DH2}) gives $s_T / \sL \propto \Rey^{1/2}$ \citep{Damkohler1940,Chaudhuri2012}. The difference lies in the mixing length: $\ell_0$ in that estimate, $\dL$ in (\ref{eq:DD}).

(iii) \textit{Sub-flame scales.} The part of (\ref{eq:S}) with $\kh > 1$, for which $\hat f_L \approx 1$ when $\ell_0 \gg \dL$, is
\begin{align}
  \mathcal{S}^{>} & = \, \Cs(\beta / 2)^{1/3} \Ka^{1/2} \left[\Gamma(-1/3, \, \beta \Ka^{-1/2} / 2) - \Gamma(-1/3, \, \beta / 2)\right] \nonumber \\
  & = \, 3 \Cs \Ka^{2/3} - c_\beta \Cs \Ka^{1/2} + \mathcal{O}(\Ka^{1/6}),
  \label{eq:Sgt}
\end{align}
which retains the leading-order $\Ka^{2/3}$ scaling of \citet{Nivarti2019} and hence their $(s_T / \sL) / (A_T / A) \sim \Ka^{1/3}$, now with a dimensionally consistent integrand. In the inertial range, the remaining large-scale part is $\mathcal{S}^{<} \approx 3 \Cs(\up / \sL - \Ka^{2/3})$, so that $\mathcal{S}^{<} / \mathcal{S}^{>} \approx (\ell_0 / \dL)^{1/3} - 1$. Eddies larger than the flame therefore dominate once $\ell_0 / \dL \gtrsim 8$. The total, however, is a smooth function of $\up / \sL$ and $\Rey$ in which $\dL$ plays no distinct role, so a division of turbulence into `large' and `small' scales at $\dL$ is not needed to predict $s_T$.

\subsection{Equivalence with power-law flame wrinkling}\label{sec:wrinkling}

\citet{Charlette2002a} modelled the unresolved wrinkling of a flame in a large-eddy simulation with filter width $\Delta$ as a power law between an outer cut-off $\Delta$ and an inner cut-off $\eta_c$,
\begin{equation}
  \Xi_\Delta = \left(1 + \frac{\Delta}{\eta_c}\right)^{\beta_w}, \qquad \frac{\Delta}{\eta_c} = \Gamma_\Delta \, \frac{\up_\Delta}{\sL}
  \quad \longrightarrow \quad
  \Xi_\Delta = \left(1 + \Gamma_\Delta \frac{\up_\Delta}{\sL}\right)^{\beta_w}.
  \label{eq:charlette}
\end{equation}
Here $\eta_c$ follows from an equilibrium between the production of surface by sub-filter strain, $\Gamma_\Delta \up_\Delta / \Delta$, and its destruction by propagation, $\sL / \eta_c$. The efficiency function $\Gamma_\Delta$ is a spectral superposition of the straining by all eddies smaller than $\Delta$ \citep{MeneveauPoinsot1991,Colin2000}. Evaluated at $\Delta = \ell_0$, the limit noted by \citet{Charlette2002a}, and with $\beta_w = 1/2$, (\ref{eq:charlette}) is identical to (\ref{eq:main}) with $\Gamma_{\ell_0} \equiv \Cs \, \mathcal{F}(\Rey)$.

Equation (\ref{eq:main}) can therefore be read in two ways: as Damk\"ohler's second hypothesis with a spectral eddy diffusivity, or as his first hypothesis with a power-law wrinkling of exponent $1/2$ and a spectrally derived efficiency. The burning velocity cannot distinguish between them. \citet{Charlette2002a} chose $\beta_w = 0.5$ from comparisons with measured $s_T$, and their dynamic procedure \citep{Charlette2002b} returned local exponents concentrated between 0.5 and 0.6 (supplementary figure~S1). They also argued that the fractal value $\beta_w = 1/3$, required to recover Damk\"ohler's large-$\up$ limit $s_T \propto \up$, is unlikely to be reached in practice. The square root in (\ref{eq:sT}) is thus consistent with the observed behaviour of large-scale wrinkling, which resolves the second difficulty raised in \S\,\ref{sec:intro}. The equivalence concerns the consumption rate rather than the area itself. The sub-$\dL$ contribution (\ref{eq:Sgt}) acts through transport within the flame \citep{Nivarti2019} and is not expected to appear in a measured $A_T / A$.

%=====================================================================
\section{Comparison with experiments}\label{sec:comparison}

\begin{table}
  \begin{center}
  \def~{\hphantom{0}}
  \begin{tabular}{clllcccc}
    & Dataset & Configuration & $\phi$ [-] & $\up / \sL$ [-] & $\ell_0 / \dL$ [-] & $\overline{\Lambda}$ [-]& $n$ \\[3pt]
    \sym{michigan2017}  & Michigan \citep{Wabel2017}        & piloted Bunsen, 21.6\,mm & 0.75    & 5.2--163  & 17--114 & 0.5  & 18 \\
    \sym{michigan2005}  & Michigan \citep{Filatyev2005}     & piloted slot, 25.4\,mm   & 1.0     & 0.68--4.8 & 8.9--37 & 0.5  & ~7 \\
    \sym{toronto2013}   & Toronto \citep{YuenGulder2013}    & piloted Bunsen, 11.2\,mm & 0.6--1.0 & 3.2--24  & 1.6--3.8 & 0.5 & ~8 \\
    \sym{toronto2014}   & Toronto \citepalias{Tamadonfar2014} & piloted Bunsen, 11.1\,mm & 0.7--1.0 & 1.8--12  & 13--40  & 0.5  & 16 \\
    \sym{tohoku1996}    & Tohoku \citep{Kobayashi1996}      & piloted Bunsen, 20\,mm   & 0.9     & 0.29--1.4 & 19--42  & 0.5  & ~4 \\
    \sym{aachen2000}    & Aachen \citep{Plessing2000}       & low-swirl, 50.8\,mm      & 0.7     & 3.5--19   & 43      & 0.5  & ~6 \\
    \sym{ottawa1995}    & Ottawa \citep{Smallwood1995}      & Bunsen, 25\,mm           & 0.7--1.0 & 0.62--2.4 & 33--141 & 0.05 & ~6 \\
    \sym{damkohler1940} & \citet{Damkohler1940}             & Bunsen, 1.4--2.7\,mm      & --      & 0--5.8    & 11--21  & --   & 15 \\
  \end{tabular}
  \caption{Datasets, with the symbols used in figures~\ref{fig:data}, \ref{fig:comparison} and \ref{fig:landscape}. All are methane--air flames at 0.1\,MPa and about 300\,K except those of \citet{Damkohler1940} (propane--butane--oxygen). Here $\phi$ is the equivalence ratio, $\overline{\Lambda}$ the mean progress variable of the isosurface defining $s_T$, and $n$ the number of points.}
  \label{tab:data}
  \end{center}
\end{table}
% TODO (data, before submission):
%  (1) delta_L is defined differently across the source papers (thermal thickness for Michigan 2017/Toronto, nu/s_L
%      for Tohoku, D(1500 K)/s_L for Aachen). The ratio of reported Re to (u'/s_L)(l0/dL) is ~5.5 (Michigan 2017),
%      ~6.6 (Toronto 2013), ~8.9 (Michigan 2005). Harmonise to delta_L = nu/s_L via RE_CORR in make_figures.ipynb.
%      Predictions depend on delta_L only through F(Re), so s_T changes by <~20%, but Re and Ka on the axes move.
%  (2) confirm burner type / pilot for Tohoku and Michigan 2005, Damkohler burner diameters (xlsx gives 1.385-2.72).
%  (3) confirm that wf_max/wf_min in Damkohler_1947.xlsx is s_T/s_L and how u'/s_L, l0/dL were estimated.
%  (4) notebook figures used c = 0.05 for Toronto (2013) in some plots and c = 0.5 in others; c = 0.5 used here.

Table~\ref{tab:data} lists the experiments, which span $0.02 < \Ka < 310$ and $10 < \Rey < 2 \times 10^4$ (figure~\ref{fig:data}\textit{a}). For uniformity, $s_T$ is the global consumption speed on the $\overline{\Lambda} = 0.5$ isosurface wherever the source reports it, and $\Rey$ and $\Ka$ are evaluated from the tabulated $\up / \sL$ and $\ell_0 / \dL$ using (\ref{eq:groups}). The spectral constants are those of non-reacting turbulence: $C = 1.5$, $\beta = 5.2$ and $c_L = 6.78$ \citep{Pope2000}, with $\De(1) = 1.88$ from \citet{Chollet1984}, so that $\Cs = 2.64$. We take $p_0 = 0$, a spectrum that is flat at low wavenumbers, which gave the best agreement in the comparisons of \citet{Walker2019}; $p_0 = 2$ and $4$ lower the predictions by medians of 12\,\% and 20\,\% over the present data. Apart from this choice, no constant is adjusted to the data.
% TODO: confirm the stated reason for p0 = 0 (thesis section 4.2 / draft notes) and whether it was chosen per dataset.

Figure~\ref{fig:comparison}(\textit{b}) compares (\ref{eq:main}) with the measurements. Over a more than twenty-fold range of $s_T / \sL$, 71\,\% of the 80 points lie within a factor of 1.5 of the prediction, and the r.m.s. logarithmic deviation is 0.38. The two high-$\Ka$ datasets are captured best. For Michigan \citep{Wabel2017}, at $3 < \Ka < 310$ and $\Rey$ up to $2 \times 10^4$, the geometric-mean ratio of measurement to prediction is 1.06 with an r.m.s.\ logarithmic deviation of 0.21; Toronto \citep{YuenGulder2013} has the same mean ratio with a deviation of 0.20. These are the conditions under which $A_T / A$ (blue symbols) falls furthest below $s_T / \sL$, and the prediction tracks $s_T$ rather than the area. The remaining datasets show systematic offsets: the low-$\Ka$ Bunsen flames of Tohoku, Ottawa and Michigan (2005) lie above the prediction by factors of 1.4--1.8, the low-swirl flames of Aachen by a factor of 2, and the Toronto data of \citet{Tamadonfar2014} at 0.8 times the prediction.
% TODO: numbers above are from make_figures.ipynb (printed summary); re-run after the delta_L harmonisation.
Offsets of this size are comparable to the effect of the isosurface choice alone, which changes the reported $s_T$ by up to a factor of two in the Toronto data \citep{YuenGulder2013}. They are also comparable to the known sensitivity of $s_T$ to flame configuration \citep{Driscoll2008,VermaLipatnikov2016}. There is, however, a weak trend with $\Ka$: points with $\Ka \leqslant 1$ lie above the prediction by a factor of 1.3 on average, whereas those with $\Ka > 1$ are unbiased (mean ratios of 1.00 for $1 < \Ka \leqslant 10$ and 1.07 for $\Ka > 10$). One candidate for the low-$\Ka$ excess is hydrodynamic (Darrieus--Landau) instability, which adds wrinkling when $\up / \sL$ is small \citep{Chaudhuri2011} and is absent from the present formulation.

\begin{figure}
  \centering
  \begin{minipage}[t]{0.48\textwidth}\raggedright (\textit{a})\\[-2pt]
    \includegraphics[width=\linewidth]{fig4a_bending}\end{minipage}\hfill
  \begin{minipage}[t]{0.52\textwidth}\raggedright (\textit{b})\\[-2pt]
    \includegraphics[width=\linewidth]{fig4b_landscape}\end{minipage}
  \caption{(\textit{a}) Measured $s_T / \sL$ against $\up / \sL$, with (\ref{eq:main}) at three Reynolds numbers; thick grey lines show the weak- and strong-turbulence limits at $\Rey = 10^3$. (\textit{b}) Contours of (\ref{eq:main}) in the $(\Rey, \Ka)$ plane, labelled with the predicted $s_T / \sL$; symbols are coloured by the measured $s_T / \sL$ on the same scale. Symbols as in table~\ref{tab:data}.}
  \label{fig:landscape}
\end{figure}

Figure~\ref{fig:landscape}(\textit{a}) shows the data against $\up / \sL$. The prediction passes from the linear growth of limit (i) to the square-root growth of limit (ii) near $\up / \sL \approx 1 / (\Cs \, \mathcal{F})$, which is 0.2--0.4 for $10^2 \leqslant \Rey \leqslant 10^4$. The data follow this bending over three decades of $\up / \sL$, and the spread in $\mathcal{F}(\Rey)$ between $\Rey = 10^2$ and $10^4$ accounts for part of their scatter at fixed $\up / \sL$. Figure~\ref{fig:landscape}(\textit{b}) maps the prediction over the regime landscape. At $\Rey \gtrsim 10^3$, where $\mathcal{F}$ is nearly constant, the contours approach lines of constant $\up / \sL$, $\Ka \propto \Rey^{-1/2}$. At lower $\Rey$, the decrease of $\mathcal{F}$ steepens them, and there $s_T$ becomes sensitive to $\Rey$ at fixed $\up / \sL$. The measured values (symbol colour) increase across the plane in the same direction as the contours, from the low-$\Rey$, low-$\up / \sL$ corner to the high-$\Rey$, high-$\Ka$ Michigan flames. Supplementary figure~S2 shows the same measurements as $(s_T / \sL) / \Ka$ against $\Rey$.

%=====================================================================
\section{Conclusions}\label{sec:conclusions}

The diffusive enhancement of the burning velocity proposed by \citet{Nivarti2019} for sub-flame eddies has been extended to the whole turbulence spectrum and made dimensionally consistent. Its scale-dependent diffusivity is taken from the spectral eddy diffusivity of two-point closures. For a model spectrum, the total enhancement reduces to $\Cs \, \mathcal{F}(\Rey) \, \up / \sL$, with $\mathcal{F}$ a weak, saturating function of the Reynolds number. The resulting burning velocity, (\ref{eq:main}), recovers three established results: Damk\"ohler's linear law in weak turbulence, the bending of $s_T$ in strong turbulence, and the $\Ka^{1/3}$ excess of $s_T / \sL$ over $A_T / A$ at high Karlovitz number. With spectral constants from non-reacting turbulence, it follows eight datasets over four decades of $\Ka$ and $\Rey$, mostly to within a factor of 1.5. The same expression is the power-law wrinkling factor of \citet{Charlette2002a} with exponent $1/2$, so area enhancement by large eddies and diffusive enhancement by small eddies need not be treated as distinct mechanisms when predicting $s_T$.

Two implications follow. First, evaluated at a filter scale $\Delta$ with the sub-filter intensity $\up_\Delta$ and Reynolds number $\Rey_\Delta$, (\ref{eq:main}) provides a sub-filter wrinkling factor of the form (\ref{eq:charlette}) whose efficiency function, $\Cs \, \mathcal{F}(\Rey_\Delta)$, follows from the spectrum rather than from a fit to combustion data. Second, the remaining differences between datasets may be explained by the quantities the theory neglects: anisotropy of the approaching flow and the effect of heat release on its spectrum, differences in the definitions of $s_T$ and $\dL$, and differential diffusion, stretch and quenching. Measurements of the turbulence spectrum upstream of the flame, and simulations at matching $\up / \sL$ and $\Rey$ in different configurations, would test whether the dependence on these two groups alone is complete.

\backsection[Supplementary data]{\label{SupMat}Supplementary material is available at \url{https://doi.org/} \dots.}

\backsection[Acknowledgements]{The authors thank Dr Aaron Skiba for helping interpret the experimental data. The authors thank Professors Denis Veynante, Peter Davidson, Andy Aspden and Matthew Juniper for helpful discussions and critical feedback. 

This report was prepared as an account of work sponsored by an agency of the United States Government. Neither the U.S. Government nor any agency thereof, nor any of their employees, makes any warranty, express or implied, or assumes any legal liability or responsibility for the accuracy, completeness, or usefulness of any information, apparatus, product, or process disclosed, or represents that its use would not infringe privately owned rights. Reference herein to any specific commercial product, process, or service by trade name, trademark, manufacturer, or otherwise does not necessarily constitute or imply its endorsement, recommendation, or favoring by the U.S. Government or any agency thereof. The views and opinions of authors expressed herein do not necessarily state or reflect those of the U.S. Government or any agency thereof.
}

\backsection[Funding]{This material is based upon work supported by the U.S. Department of Energy, Office of Science, Office of Advanced Scientific Computing Research, Department of Energy Computational Science Graduate Fellowship under Award Number DE-SC0024386. M.D.W. was further supported by the Churchill Scholarship from the Winston Churchill Foundation of the United States. G.V.N. acknowledges the funding support of EPSRC grant EP/P022286/1.}

\backsection[Declaration of interests]{%A Competing Interests statement is now mandatory in the manuscript PDF. Please note that if there are no conflicts of interest, the declaration in your PDF should read as follows: {\bf Declaration of Interests}.
The authors report no conflict of interest.}

\backsection[Data availability statement]{The data that support the findings of this study are openly available on GitHub at \url{https://github.com/mw6136/spectral}.}

\backsection[Author ORCIDs]{ M. D. Walker, https://orcid.org/0000-0002-4765-6539; G. V. Nivarti, https://orcid.org/0000-0002-8326-1147; M. E. Mueller, https://orcid.org/0000-0002-5699-1275; R. S. Cant, https://orcid.org/0000-0003-1851-6665; S. Hochgreb, https://orcid.org/0000-0001-7192-4786}

\backsection[Author contributions]{M. D. Walker: Writing -- original draft, Validation, Software, Methodology, Investigation, Data Curation. G. V. Nivarti: Methodology, Investigation, Conceptualization, Writing -- review \& editing. M. E. Mueller: Writing -- review \& editing, Supervision. R. S. Cant: Methodology, Investigation, Writing -- review \& editing, Supervision. S. Hochgreb: Conceptualization, Writing -- review \& editing, Supervision.}

\bibliographystyle{jfm}
\bibliography{jfm}


\end{document}